% =====================================================
\textbf{UC-FIR-01: Wysłanie prośby o zadatek}

\begin{longtable}{|p{4cm}|p{11cm}|}
\hline
\textbf{Element} & \textbf{Opis} \\
\hline

Identyfikator & UC-FIR-01 \\
\hline

Nazwa & Wysłanie prośby o zadatek \\
\hline

Opis & 
Przygotowanie danych do przelewu i powiadomienie klienta o konieczności wpłaty zadatku w przypadku braku pojazdu na placu. \\
\hline

Aktor główny & System (Kontekst Fakturowania) \\
\hline

Aktorzy pomocniczy & Klient \\
\hline

Warunki wstępne & 
Kontekst otrzymał zdarzenie \textit{VehicleIsNotOnStock}. \\
\hline

Warunki końcowe & 
Klient został poinformowany o wymaganej wpłacie, emitowane zdarzenie \textit{AdvancePaymentRequested}. \\
\hline

Scenariusz główny & 
\begin{enumerate}
    \item Kontekst fakturowania reaguje na zdarzenie \textit{VehicleIsNotOnStock}.
    \item System generuje numer rachunku bankowego oraz kwotę zadatku dla zamówienia.
    \item System wysyła do klienta informację (np. e-mail) z danymi do przelewu.
    \item Kontekst emituje zdarzenie \textit{AdvancePaymentRequested}.
\end{enumerate} \\
\hline

Scenariusze alternatywne & 
\begin{itemize}
    \item [A1] Błąd danych: Brak wymaganych informacji do wygenerowania prośby, system emituje \textit{ErrorDuringPaymentRequest}.
\end{itemize} \\
\hline
\end{longtable}

% =====================================================
\textbf{UC-FIR-02: Stworzenie faktury końcowej}

\begin{longtable}{|p{4cm}|p{11cm}|}
\hline
\textbf{Element} & \textbf{Opis} \\
\hline

Identyfikator & UC-FIR-02 \\
\hline

Nazwa & Stworzenie faktury końcowej \\
\hline

Opis & 
Automatyczne wystawienie faktury końcowej po potwierdzeniu rezerwacji pojazdu z placu. \\
\hline

Aktor główny & System (Kontekst Fakturowania) \\
\hline

Aktorzy pomocniczy & - \\
\hline

Warunki wstępne & 
Kontekst otrzymał zdarzenie \textit{VehicleReservedFromStock} lub {VehicleDeliveredToStock}. \\
\hline

Warunki końcowe & 
Faktura PDF utworzona i przypisana do zamówienia, emitowane zdarzenie \textit{InvoiceCreated}. \\
\hline

Scenariusz główny & 
\begin{enumerate}
    \item Kontekst fakturowania reaguje na zdarzenie \textit{VehicleReservedFromStock}.
    \item System tworzy fakturę na kwotę pozostałą do zapłaty (uwzględniając ewentualne wpłaty zadatku).
    \item System wysyła do klienta informację (np. e-mail) z danymi do przelewu.
    \item System generuje plik PDF faktury i przypisuje go do zamówienia.
    \item Kontekst emituje zdarzenie \textit{InvoiceCreated}.
\end{enumerate} \\
\hline

Scenariusze alternatywne & 
\begin{itemize}
    \item [A1] Błąd generowania dokumentu: System nie może utworzyć pliku, emituje \textit{ErrorDuringInvoiceCreation}.
\end{itemize} \\
\hline
\end{longtable}

% =====================================================
\textbf{UC-FIR-03: Procesowanie płatności}

\begin{longtable}{|p{4cm}|p{11cm}|}
\hline
\textbf{Element} & \textbf{Opis} \\
\hline

Identyfikator & UC-FIR-03 \\
\hline

Nazwa & Procesowanie płatności \\
\hline

Opis & 
Zaksięgowanie przelewu bankowego i rozliczenie salda zamówienia. \\
\hline

Aktor główny & Księgowy \\
\hline

Aktorzy pomocniczy & System fakturowania \\
\hline

Warunki wstępne & 
W lokalnej bazie istnieje otwarta należność, księgowy zidentyfikował wpłatę na wyciągu. \\
\hline

Warunki końcowe & 
Płatność zarejestrowana, saldo zamówienia zaktualizowane, emitowane \textit{PaymentRegistered} (ewent. \textit{SettlementCompleted}). \\
\hline

Scenariusz główny & 
\begin{enumerate}
    \item Księgowy importuje plik wyciągu bankowego do systemu.
    \item System paruje przelew z zamówieniem (po tytule/ID). W przypadku braku dopasowania automatycznego, księgowy dokonuje weryfikacji ręcznej.
    \item System aktualizuje status płatności na 'REGISTERED'.
    \item Kontekst emituje zdarzenie \textit{PaymentRegistered}.
    \item Jeżeli kwota wpłaty jest równa kwocie zadatku system emituje zdarzenie \textit{AdvancePaymentRegistered}
    \item System weryfikuje sumę wpłat. Jeśli saldo = 0 PLN, kontekst emituje \textit{SettlementCompleted}.
\end{enumerate} \\
\hline

Scenariusze alternatywne & 
\begin{itemize}
    \item [A1] Błędna kwota: Kwota wpłaty różni się od wymaganej, system księguje częściowo, status płatności to \textit{PARTIAL\_PAYMENT}.
\end{itemize} \\
\hline
\end{longtable}


% =====================================================
\textbf{UC-FIR-04: Monitorowanie terminów płatności}

\begin{longtable}{|p{4cm}|p{11cm}|}
\hline
\textbf{Element} & \textbf{Opis} \\
\hline

Identyfikator & UC-FIR-04 \\
\hline

Nazwa & Monitorowanie terminów płatności \\
\hline

Opis & 
Zautomatyzowany proces cyklicznie weryfikujący terminy płatności wystawionych faktur. Jeśli płatność jest przeterminowana o 30 dni, system uruchamia procedurę zwalniania rezerwacji. \\
\hline

Aktor główny & System (Kontekst Fakturowania) \\
\hline

Aktorzy pomocniczy & - \\
\hline

Warunki wstępne & 
W systemie istnieją wygenerowane faktury końcowe (po zdarzeniu \textit{InvoiceCreated}), dla których saldo jest różne od zera. \\
\hline

Warunki końcowe & 
Przeterminowane faktury otrzymują odpowiedni status, a system emituje zdarzenie \textit{PaymentDeadlineExpired} dla każdego nieopłaconego w terminie zamówienia. \\
\hline

Scenariusz główny & 
\begin{enumerate}
    \item Harmonogram zadań (np. raz na dobę w nocy) uruchamia proces weryfikacji.
    \item System wyszukuje wszystkie otwarte faktury ze statusem oczekującym na płatność.
    \item System weryfikuje datę wystawienia dla każdej z tych faktur.
    \item System identyfikuje zamówienia, dla których od daty wystawienia faktury upłynęło 30 dni.
    \item System zmienia wewnętrzny status tych faktur na "Przeterminowana".
    \item Kontekst emituje zdarzenie \textit{PaymentDeadlineExpired} na szynę danych.
\end{enumerate} \\
\hline

Scenariusze alternatywne & 
\begin{itemize}
    \item [A1] Brak przeterminowanych faktur: W kroku 4 system nie znajduje żadnej faktury starszej niż 30 dni. Proces kończy się bez żadnych zmian i bez emisji zdarzeń.
\end{itemize} \\
\hline
\end{longtable}